\olchapter{fol}{syn}{Syntax of First-Order Logic}



\olfileid{fol}{syn}{itx}

\olsection{Introduction}

In order to develop the theory and metatheory of first-order logic, we
must first define the syntax and semantics of its expressions.  The
expressions of first-order logic are terms and !!{formula}s.  Terms
are formed from !!{variable}s, !!{constant}s, and !!{function}s.
!!^{formula}s, in turn, are formed from !!{predicate}s together with
terms (these form the smallest, ``atomic'' !!{formula}s), and then
from atomic !!{formula}s we can form more complex ones using logical
connectives and quantifiers.  There are many different ways to set
down the formation rules; we give just one possible one. Other systems
will chose different symbols, will select different sets of
connectives as primitive, will use parentheses differently (or even not
at all, as in the case of so-called Polish notation).  What all
approaches have in common, though, is that the formation rules define
the set of terms and !!{formula}s \emph{inductively}. If done
properly, every expression can result essentially in only one way
according to the formation rules.  The inductive definition resulting
in expressions that are \emph{uniquely readable} means we can give
meanings to these expressions using the same method---inductive
definition.





\olfileid{fol}{syn}{fol}

\olsection{First-Order Languages}


Expressions of first-order logic are built up from a basic vocabulary
containing \emph{!!{variable}s}, \emph{!!{constant}s},
\emph{!!{predicate}s} and sometimes \emph{!!{function}s}.  From them,
together with logical connectives, quantifiers, and punctuation
symbols such as parentheses and commas, \emph{terms} and
\emph{!!{formula}s} are formed.

\begin{explain}
Informally, !!{predicate}s are names for properties and relations,
!!{constant}s are names for individual objects, and !!{function}s are
names for mappings.  These, except for the !!{identity}~$\eq$, are the
\emph{non-logical symbols} and together make up a language.  Any
first-order language~$\Lang L$ is determined by its non-logical
symbols.  In the most general case, $\Lang L$ contains infinitely many
symbols of each kind.
\end{explain}

In the general case, we make use of the following symbols in
first-order logic:

\begin{enumerate}
\item Logical symbols
\begin{enumerate}
\item Logical connectives:
  \startycommalist
  \ycomma $\lnot$ (negation)
  \ycomma $\land$ (conjunction)
  \ycomma $\lor$ (disjunction)
  \ycomma $\lif$ (!!{conditional})
  
  \ycomma $\lforall$ (universal quantifier)
  \ycomma $\lexists$ (existential quantifier).
The propositional constant for !!{falsity}~$\lfalse$.

\item The two-place !!{identity}~$\eq$.
\item A !!{denumerable}s set of !!{variable}s: $\Obj v_0$, $\Obj v_1$, $\Obj
  v_2$, \dots
\end{enumerate}
\item Non-logical symbols, making up the \emph{standard
  language} of first-order logic
\begin{enumerate}
\item A !!{denumerable}s set of $n$-place !!{predicate}s for each $n>0$: $\Obj
  A^n_0$, $\Obj A^n_1$, $\Obj A^n_2$, \dots
\item A !!{denumerable}s set of !!{constant}s: $\Obj c_0$, $\Obj c_1$, $\Obj
  c_2$, \dots.
\item A !!{denumerable}s set of $n$-place !!{function}s for each $n>0$:
  $\Obj f^n_0$, $\Obj f^n_1$, $\Obj f^n_2$, \dots
\end{enumerate}
\item Punctuation marks: (, ), and the comma.
\end{enumerate}

Most of our definitions and results will be formulated for the full
standard language of first-order logic.  However, depending on the
application, we may also restrict the language to only a few
!!{predicate}s, !!{constant}s, and !!{function}s.

\begin{ex}
The language~$\Lang L_A$ of arithmetic contains a single two-place
!!{predicate}~$<$, a single !!{constant}~$\Obj 0$, one one-place
!!{function}~$\prime$, and two two-place !!{function}s~$+$ and~$\times$.
\end{ex}

\begin{ex}
The language of set theory~$\Lang L_Z$ contains only the single
two-place !!{predicate}~$\in$.
\end{ex}

\begin{ex}
The language of orders~$\Lang L_\le$ contains only the two-place
!!{predicate}~$\le$.
\end{ex}

Again, these are conventions: officially, these are just aliases,
e.g., $<$, $\in$, and $\le$ are aliases for $\Obj A^2_0$, $\Obj 0$ for
$\Obj c_0$, $\prime$ for $\Obj f^1_0$, $+$ for $\Obj f^2_0$, $\times$ for
$\Obj f^2_1$.


In addition to the primitive connectives and
quantifiers introduced
above, we also use the following \emph{defined} symbols:
\startycommalist
  
  
  
  
  \ycomma $\liff$ (!!{biconditional})
  
  
  
  \ycomma !!{truth}~$\ltrue$.


\begin{explain}
A defined symbol is not officially part of the language, but is
introduced as an informal abbreviation: it allows us to abbreviate
formulas which would, if we only used primitive symbols, get quite
long.  This is obviously an advantage.  The bigger advantage, however,
is that proofs become shorter.  If a symbol is primitive, it has to be
treated separately in proofs. The more primitive symbols, therefore,
the longer our proofs.
\end{explain}




\begin{intro}
You may be familiar with different terminology and symbols than the
ones we use above. Logic texts (and teachers) commonly use 
$\sim$, $\neg$, or~!\ for ``negation'', $\wedge$, $\cdot$, or $\&$
for ``conjunction''.  Commonly used symbols for the ``conditional'' or
``implication'' are $\rightarrow$, $\Rightarrow$, and $\supset$.
Symbols for ``biconditional,'' ``bi-implication,''
  or ``(material) equivalence'' are $\leftrightarrow$,
  $\Leftrightarrow$, and $\equiv$.
The $\lfalse$ symbol is variously called
  ``falsity,'' ``falsum,'', ``absurdity,'' or ``bottom.''
The $\ltrue$ symbol is variously called
  ``truth,'' ``verum,'' or ``top.''

It is conventional to use lower case letters (e.g., $a$, $b$, $c$) from
the beginning of the Latin alphabet for !!{constant}s (sometimes called
names), and lower case letters from the end (e.g., $x$, $y$, $z$) for
!!{variable}s. Quantifiers combine with !!{variable}s, e.g., $x$;
notational variations include $\forall x$, $(\forall x)$, $(x)$, $\Pi x$,
$\bigwedge_x$ for the universal quantifier and $\exists x$, $(\exists
x)$, $(Ex)$, $\Sigma x$, $\bigvee_x$ for the existential quantifier.
\end{intro}

\begin{explain}
We might treat all the propositional operators and both quantifiers as
primitive symbols of the language.  We might instead choose a smaller
stock of primitive symbols and treat the other !!{operator}s as
defined. ``Truth functionally complete'' sets of Boolean operators
include $\{ \lnot, \lor \}$, $\{ \lnot, \land \}$, and $\{ \lnot,
\lif\}$---these can be combined with either quantifier for an
expressively complete first-order language.

You may be familiar with two other !!{operator}s: the Sheffer
stroke~$|$ (named after Henry Sheffer), and Peirce's
arrow~$\downarrow$, also known as Quine's dagger.  When given their
usual readings of ``nand'' and ``nor'' (respectively), these operators
are truth functionally complete by themselves.
\end{explain}





\olfileid{fol}{syn}{frm}

\olsection{Terms and \printtoken{P}{formula}}

Once a first-order language~$\Lang L$ is given, we can define
expressions built up from the basic vocabulary of~$\Lang L$.  These
include in particular \emph{terms} and \emph{!!{formula}s}.

\begin{defn}[Terms]
\ollabel{defn:terms}
The set of \emph{terms}~$\Trm[L]$ of~$\Lang L$ is
defined inductively by:
\begin{enumerate}
\item Every !!{variable} is a term.
\item Every !!{constant} of~$\Lang L$ is a term.
\item If $f$ is an $n$-place !!{function} and $t_1$, \dots, $t_n$
  are terms, then $\Atom{f}{t_1, \ldots, t_n}$ is a term.

  Nothing else is a term.
\end{enumerate}
A term containing no !!{variable}s is a \emph{closed term}.
\end{defn}

\begin{explain}
The !!{constant}s appear in our specification of the language and the
terms as a separate category of symbols, but they could instead have
been included as zero-place !!{function}s.  We could then do without
the second clause in the definition of terms. We just have to
understand $\Atom{f}{t_1, \ldots, t_n}$ as just $f$ by itself if $n =
0$.
\end{explain}

\begin{defn}[Formulas]
\ollabel{defn:formulas}
The set of \emph{!!{formula}s}~$\Frm[L]$ of the language~$\Lang L$
is defined inductively as follows:
\begin{enumerate}
$\lfalse$ is an atomic !!{formula}.



\item If $R$ is an $n$-place !!{predicate} of~$\Lang L$ and $t_1$, \dots,
  $t_n$ are terms of~$\Lang L$, then $\Atom{R}{t_1,\ldots, t_n}$ is an
  atomic !!{formula}.

\item If $t_1$ and $t_2$ are terms of~$\Lang L$, then $\Atom{\eq}{t_1, t_2}$
  is an atomic !!{formula}.

If $!A$ is !!a{formula}, then $\lnot !A$ is
  !!a{formula}.

If $!A$ and $!B$ are !!{formula}s, then $(!A \land
  !B)$ is !!a{formula}.

If $!A$ and $!B$ are !!{formula}s, then $(!A \lor !B)$
  is !!a{formula}.

If $!A$ and $!B$ are !!{formula}s, then $(!A \lif !B)$
  is !!a{formula}.



If $!A$ is !!a{formula} and $x$ is !!a{variable},
  then $\lforall[x][!A]$ is !!a{formula}.

If $!A$ is !!a{formula} and $x$ is !!a{variable},
  then $\lexists[x][!A]$ is !!a{formula}.

Nothing else is !!a{formula}.
\end{enumerate}
\end{defn}

\begin{explain}
The definitions of the set of terms and that of !!{formula}s are
\emph{inductive definitions}.  Essentially, we construct the set of
!!{formula}s in infinitely many stages.  In the initial stage, we
pronounce all atomic formulas to be formulas; this corresponds to the
first few cases of the definition, i.e., the cases for

$\lfalse$, 
$\Atom{R}{t_1,\dots,t_n}$ and $\Atom{\eq}{t_1,t_2}$. ``Atomic
!!{formula}'' thus means any !!{formula} of this form.

The other cases of the definition give rules for constructing new
!!{formula}s out of !!{formula}s already constructed.  At the second
stage, we can use them to construct !!{formula}s out of atomic
!!{formula}s.  At the third stage, we construct new formulas from the
atomic formulas and those obtained in the second stage, and so on.  A
!!{formula} is anything that is eventually constructed at such a
stage, and nothing else.
\end{explain}

By convention, we write $\eq$ between its arguments and leave out the
parentheses: $\eq[t_1][t_2]$ is an abbreviation for
$\Atom{\eq}{t_1,t_2}$.  Moreover, $\lnot \Atom{\eq}{t_1,t_2}$ is
abbreviated as $\eq/[t_1][t_2]$. When writing a formula $(!B \ast !C)$
constructed from $!B$, $!C$ using a two-place connective~$\ast$, we
will often leave out the outermost pair of parentheses and write
simply~$!B \ast !C$.

\begin{intro}
Some logic texts require that the !!{variable}~$x$ must occur in~$!A$
in order for
$\lexists[x][!A]$ 
and 
$\lforall[x][!A]$ 
to count as
!!{formula}s.
Nothing bad happens if you don't require this, and it makes things
easier.
\end{intro}


\begin{defn}
Formulas constructed using the defined operators are to be understood
as follows:

\begin{tagenumerate}{defTrue,defFalse,defNot,defOr,defAnd,defIf,defIff,defEx,defAll}

$\ltrue$ abbreviates
  $\lnot\lfalse$.











$!A \liff !B$ abbreviates $(!A \lif !B) \land (!B
  \lif !A)$.




\end{tagenumerate}
\end{defn}


If we work in a language for a specific application, we will often
write two-place !!{predicate}s and !!{function}s between the
respective terms, e.g., $t_1 < t_2$ and $(t_1 + t_2)$ in the language
of arithmetic and $t_1 \in t_2$ in the language of set theory.  The
successor function in the language of arithmetic is even written
conventionally \emph{after} its argument:~$t'$.  Officially, however,
these are just conventional abbreviations for $\Atom{\Obj A^2_0}{t_1, t_2}$,
$\Obj f^2_0(t_1, t_2)$, $\Atom{\Obj A^2_0}{t_1, t_2}$ and $f^1_0(t)$,
respectively.

\begin{defn}[Syntactic identity]
The symbol $\ident$ expresses syntactic identity between strings of
symbols, i.e., $!A \ident !B$ iff $!A$ and $!B$ are strings of symbols
of the same length and which contain the same symbol in each place.
\end{defn}

The $\ident$ symbol may be flanked by strings obtained by
concatenation, e.g., $!A \ident (!B \lor !C)$ means: the string of
symbols~$!A$ is the same string as the one obtained by concatenating
an opening parenthesis, the string $!B$, the $\lor$ symbol, the
string~$!C$, and a closing parenthesis, in this order. If this is the
case, then we know that the first symbol of~$!A$ is an opening
parenthesis, $!A$ contains $!B$ as a substring (starting at the second
symbol), that substring is followed by~$\lor$, etc.

As terms and !!{formula}s are built up from basic elements via inductive
definitions, we can use the following induction principles to prove
things about them.

\begin{lem}[\emph{Principle of induction on terms}]
    \ollabel{lem:trmind}
    Let $\Lang L$ be a first-order language.
    If some property~$P$ is such that
    
    \begin{enumerate}
        \item it holds for every !!{variable}~$v$,
        
        \item it holds for every !!{constant}~$a$ of~$\Lang L$, and
        
        \item it holds for $f(t_1,\dotsc,t_n)$ whenever it holds for
        $t_1$,~\dots, $t_n$ and $f$~is an $n$-place
            !!{function} of~$\Lang L$
    \end{enumerate}
    (assuming $t_1$,~\dots, $t_n$ are terms of~$\Lang{L}$),
    then $P$ holds for every term in~$\Trm[L]$.
\end{lem}

\begin{prob}
    Prove \olref[fol][syn][frm]{lem:trmind}.
\end{prob}

\begin{lem}[\emph{Principle of induction on !!{formula}s}]
    \ollabel{thm:frmind}
    Let $\Lang L$ be a first-order language.
    If some property~$P$ holds for all the atomic !!{formula}s
    and is such that
    
    \begin{enumerate}
        it holds for $\lnot !A$ whenever it
            holds for~$!A$;
        it holds for $(!A \land !B)$
            whenever it holds for $!A$ and~$!B$;
        it holds for $(!A \lor !B)$
            whenever it holds for $!A$ and~$!B$;
        it holds for $(!A \lif !B)$
            whenever it holds for $!A$ and~$!B$;
        
        it holds for $\lexists[x][!A]$
            whenever it holds for~$!A$;
        it holds for $\lforall[x][!A]$
            whenever it holds for~$!A$;
  \end{enumerate}
  (assuming $!A$ and $!B$ are !!{formula}s of~$\Lang{L}$),
  then $P$ holds for all formulas in~$\Frm[L]$.
\end{lem}





\olfileid{fol}{syn}{unq}

\olsection{Unique Readability}

\begin{explain}
The way we defined !!{formula}s guarantees that every !!{formula} has
a \emph{unique reading}, i.e., there is essentially only one way of
constructing it according to our formation rules for !!{formula}s and
only one way of ``interpreting'' it.  If this were not so, we would
have ambiguous !!{formula}s, i.e., !!{formula}s that have more than
one reading or intepretation---and that is clearly something we want
to avoid.  But more importantly, without this property, most of the
definitions and proofs we are going to give will not go through.

Perhaps the best way to make this clear is to see what would happen if
we had given bad rules for forming !!{formula}s that would not
guarantee unique readability.  For instance, we could have forgotten
the parentheses in the formation rules for connectives, e.g., we might
have allowed this:
\begin{quote}
If $!A$ and $!B$ are !!{formula}s, then so is $!A \lif !B$.
\end{quote}
Starting from an atomic formula $!D$, this would allow us to form $!D
\lif !D$. From this, together with $!D$, we would get $!D \lif !D
\lif !D$.  But there are two ways to do this:
\begin{enumerate}
\item We take $!D$ to be $!A$ and $!D \lif !D$ to be $!B$.
\item We take $!A$ to be $!D \lif !D$ and $!B$ is~$!D$.  
\end{enumerate}
Correspondingly, there are two ways to
``read'' the !!{formula}~$!D \lif !D \lif !D$. It is of the form $!B
\lif !C$ where $!B$ is $!D$ and $!C$ is $!D \lif !D$, but \emph{it is
  also} of the form $!B \lif !C$ with $!B$ being $!D \lif !D$ and $!C$
being~$!D$.

If this happens, our definitions will not always work. For instance,
when we define the !!{main operator} of a formula, we say: in a
formula of the form $!B \lif !C$, the !!{main operator} is the
indicated occurrence of~$\lif$. But if we can match the formula $!D
\lif !D \lif !D$ with $!B \lif !C$ in the two different ways mentioned
above, then in one case we get the first occurrence of $\lif$ as the
!!{main operator}, and in the second case the second occurrence.  But
we intend the !!{main operator} to be a \emph{function} of the
!!{formula}, i.e., every !!{formula} must have exactly one !!{main
  operator} occurrence.
\end{explain}

\begin{lem}
The number of left and right parentheses in !!a{formula}~$!A$ are
equal.
\end{lem}

\begin{proof}
We prove this by induction on the way $!A$ is constructed. This
requires two things: (a) We have to prove first that all atomic
formulas have the property in question (the induction basis). (b) Then
we have to prove that when we construct new formulas out of given
formulas, the new formulas have the property provided the old ones do.

Let $l(!A)$ be the number of left parentheses, and $r(!A)$ the number
of right parentheses in~$!A$, and $l(t)$ and $r(t)$ similarly the
number of left and right parentheses in a term~$t$.

\begin{prob}
    Prove that for any term~$t$, $l(t) = r(t)$.
\end{prob}

\begin{enumerate}
\indcase{!A}{\lfalse}{$\indfrm$ has $0$ left and $0$
    right parentheses.}



\item \indcase{!A}{\Atom{R}{t_1,\dots,t_n}}{$l(\indfrm) = 1 + l(t_1) +
  \dots + l(t_n) = 1 + r(t_1) + \dots + r(t_n) = r(\indfrm)$. Here we
  make use of the fact, left as an exercise, that $l(t) = r(t)$ for
  any term~$t$.}

\item \indcase{!A}{\eq[t_1][t_2]}{$l(\indfrm) = l(t_1) + l(t_2) =
  r(t_1) + r(t_2) = r(\indfrm)$.}

\indcase{!A}{\lnot !B}{By induction hypothesis,
    $l(!B) = r(!B)$. Thus $l(\indfrm) = l(!B) = r(!B) =
    r(\indfrm)$.}

\item \indcase{!A}{(!B \ast !C)}{By induction hypothesis, $l(!B) =
  r(!B)$ and $l(!C) = r(!C)$.  Thus $l(\indfrm) = 1 + l(!B) + l(!C) =
  1 + r(!B) + r(!C) = r(\indfrm)$.}

\indcase{!A}{\lforall[x][!B]}{By induction
    hypothesis, $l(!B) = r(!B)$. Thus, $l(\indfrm) = l(!B) = r(!B) =
    r(\indfrm)$.}





\indcase{!A}{\lexists[x][!B]}{Similarly.}
\end{enumerate}
\end{proof}

\begin{defn}[Proper prefix]
A string of symbols $!B$ is a \emph{proper prefix} of a string of symbols~$!A$ if
concatenating $!B$ and a non-empty string of symbols yields~$!A$.
\end{defn}

\begin{lem}\ollabel{lem:no-prefix}
If $!A$ is !!a{formula}, and $!B$ is a proper prefix of $!A$, then
$!B$ is not !!a{formula}.
\end{lem}

\begin{proof}
Exercise.
\end{proof}

\begin{prob}
Prove \olref[fol][syn][unq]{lem:no-prefix}.
\end{prob}

\begin{prop}
\ollabel{prop:unique-atomic}
If $!A$ is an atomic !!{formula}, then it satisfies one, and only one
of the following conditions.
\begin{enumerate}
$!A \ident \lfalse$.

\item $!A \ident \Atom{R}{t_1,\dots,t_n}$ where $R$ is an $n$-place
  !!{predicate}, $t_1$, \dots, $t_n$ are terms, and each of $R$,
  $t_1$, \dots, $t_n$ is uniquely determined.
\item $!A \ident \eq[t_1][t_2]$ where $t_1$ and $t_2$ are uniquely
  determined terms.
\end{enumerate}
\end{prop}

\begin{proof}
Exercise.
\end{proof}

\begin{prob}
Prove \olref[fol][syn][unq]{prop:unique-atomic} (Hint: Formulate and
prove a version of \olref[fol][syn][unq]{lem:no-prefix} for terms.)
\end{prob}

\begin{prop}[Unique Readability]
Every !!{formula} satisfies one, and only one of the following conditions.
\begin{enumerate}
\item $!A$ is atomic.

$!A$ is of the form $\lnot !B$.

$!A$ is of the form $(!B \land !C)$.

$!A$ is of the form $(!B \lor !C)$.

$!A$ is of the form $(!B \lif !C)$.



$!A$ is of the form $\lforall[x][!B]$.

$!A$ is of the form $\lexists[x][!B]$.
\end{enumerate}
Moreover, in each case $!B$, or $!B$ and $!C$, are uniquely
determined.  This means that, e.g., there are no different pairs $!B$,
$!C$ and $!B'$, $!C'$ so that $!A$ is both of the form 
$(!B \lif !C)$ and $(!B' \lif !C')$.
\end{prop}

\begin{proof}
The formation rules require that if !!a{formula} is not atomic, it
must start with an opening parenthesis~(, $\lnot$,
or a quantifier. On the other hand, every !!{formula} that starts with
one of the following symbols must be atomic: !!a{predicate}, !!a{function}, !!a{constant}, $\lfalse$.

So we really only have to show that if $!A$ is of the form $(!B \ast
!C)$ and also of the form $(!B' \mathbin{\ast'} !C')$, then $!B \ident
!B'$, $!C \ident !C'$, and $\ast = {\ast'}$.

So suppose both $!A \ident (!B \ast !C)$ and $!A \ident (!B'
\mathbin{\ast'} !C')$.  Then either $!B \ident !B'$ or not.  If it is,
clearly $\ast = {\ast'}$ and $!C \ident !C'$, since they then are
substrings of $!A$ that begin in the same place and are of the same
length.  The other case is $!B \not\ident !B'$.  Since $!B$ and
$!B'$ are both substrings of $!A$ that begin at the same place, one
must be a proper prefix of the other.  But this is impossible by
\olref{lem:no-prefix}.
\end{proof}






\olfileid{fol}{syn}{mai}

\olsection{\printtoken{S}{main operator} of a Formula}

\begin{explain}
It is often useful to talk about the last operator used in
constructing !!a{formula}~$!A$.  This operator is called the \emph{main
  operator} of~$!A$. Intuitively, it is the ``outermost'' operator
of $!A$. For example, the main operator of $\lnot !A$ is $\lnot$,
the main operator of $(!A \lor !B)$ is $\lor$, etc.
\end{explain}


\begin{defn}[!!^{main operator}]
\ollabel{def:main-op}
The \emph{!!{main operator}} of !!a{formula}~$!A$ is
defined as follows:
\begin{enumerate}
\item \indcase*{!A}{!A}{$\indfrm$ has no !!{main operator}.}

\indcase{!A}{\lnot !B}{the !!{main operator} of $\indfrm$
  is~$\lnot$.}

\indcase{!A}{(!B \land !C)}{the !!{main operator} of
  $\indfrm$ is~$\land$.}

\indcase{!A}{(!B \lor !C)}{the !!{main operator} of
  $\indfrm$ is~$\lor$.}

\indcase{!A}{(!B \lif !C)}{the !!{main operator} of
  $\indfrm$ is~$\lif$.}



\indcase{!A}{\lforall[x][!B]}{the !!{main operator}
  of $\indfrm$ is~$\lforall$.}

\indcase{!A}{\lexists[x][!B]}{the !!{main operator} of
  $\indfrm$ is~$\lexists$.}
\end{enumerate}
\end{defn}

In each case, we intend the specific indicated \emph{occurrence} of
the !!{main operator} in the formula. For instance, since the formula
$((!D \lif !E) \lif (!E \lif !D))$ is of the form $(!B \lif !C)$ where
$!B$ is $(!D \lif !E)$ and $!C$ is $(!E \lif !D)$, the second
occurrence of $\lif$ is the !!{main operator}.

\begin{explain}
This is a \emph{recursive} definition of a function which maps all
non-atomic !!{formula}s to their !!{main operator} occurrence. Because
of the way !!{formula}s are defined inductively, every
!!{formula}~$!A$ satisfies one of the cases in \olref{def:main-op}.
This guarantees that for each non-atomic !!{formula}~$!A$ !!a{main
  operator} exists. Because each !!{formula} satisfies only one of these
conditions, and because the smaller !!{formula}s from which $!A$ is
constructed are uniquely determined in each case, the !!{main
  operator} occurrence of~$!A$ is unique, and so we have defined a
function.
\end{explain}

We call !!{formula}s by the names in \olref{tab:main-op} depending on
which symbol their !!{main operator}
is. Recall, however, that defined operators do not officially appear in
!!{formula}s. They are just abbreviations, so officially they cannot
be the main operator of a formula. In proofs about all !!{formula}s
they therefore do not have to be treated separately.{table}[!h]
\centering
\begin{tabular}{c | c | c}
!!^{main operator} & Type of !!{formula} & Example\\
\hline
none & atomic (!!{formula}) &
$\lfalse$,

$\Atom{R}{t_1, \dots, t_n}$,
$\eq[t_1][t_2]$\\
$\lnot$ & negation & $\lnot !A$ \\
$\land$ & conjunction & $(!A \land !B$) \\
$\lor$ & disjunction & $(!A \lor !B$) \\
$\lif$ & !!{conditional} & $(!A \lif !B$) \\
$\liff$ & !!{biconditional} & $(!A \liff !B)$ \\
$\lforall[][]$ & universal (!!{formula})& $\lforall[x][!A]$ \\
$\lexists[][]$ & existential (!!{formula})& $\lexists[x][!A]$
\end{tabular}
\caption{Main operator and names of !!{formula}s}
\ollabel{tab:main-op}
\end{table}





\olfileid{fol}{syn}{sbf}

\olsection{\printtoken{P}{subformula}}

\begin{explain}
It is often useful to talk about the !!{formula}s that ``make up'' a
given !!{formula}.  We call these its \emph{!!{subformula}s}.  Any
!!{formula} counts as !!a{subformula} of itself; a subformula of $!A$
other than $!A$ itself is a \emph{proper !!{subformula}}.
\end{explain}

\begin{defn}[Immediate !!^{subformula}]
If $!A$ is !!a{formula}, the \emph{immediate !!{subformula}s}
of $!A$ are defined inductively as follows:
\begin{enumerate}
\item Atomic !!{formula}s have no immediate !!{subformula}s.

\indcase{!A}{\lnot !B}{The only immediate
    !!{subformula} of $\indfrm$ is~$!B$.}

\item \indcase{!A}{(!B \ast !C)}{The immediate !!{subformula}s of
  $\indfrm$ are $!B$ and $!C$ ($\ast$ is any one of the two-place
  connectives).}

\indcase{!A}{\lforall[x][!B]}{The only immediate
    !!{subformula} of $\indfrm$ is~$!B$.}

\indcase{!A}{\lexists[x][!B]}{The only immediate
    !!{subformula} of $\indfrm$ is~$!B$.}
\end{enumerate}
\end{defn}

\begin{defn}[Proper !!^{subformula}]
If $!A$ is !!a{formula}, the \emph{proper !!{subformula}s}
of $!A$ are defined recursively as follows:
\begin{enumerate}
\item Atomic !!{formula}s have no proper !!{subformula}s.

\indcase{!A}{\lnot !B}{The proper !!{subformula}s of
    $\indfrm$ are~$!B$ together with all proper !!{subformula}s
    of~$!B$.}

\item \indcase{!A}{(!B \ast !C)}{The proper !!{subformula}s of
  $\indfrm$ are $!B$, $!C$, together with all proper !!{subformula}s
  of $!B$ and those of~$!C$.}

\indcase{!A}{\lforall[x][!B]}{The proper
    !!{subformula}s of $\indfrm$ are~$!B$ together with all proper
    !!{subformula}s of~$!B$.}

\indcase{!A}{\lexists[x][!B]}{The proper
    !!{subformula}s of $\indfrm$ are~$!B$ together with all proper
    !!{subformula}s of~$!B$.}
\end{enumerate}
\end{defn}

\begin{defn}[!!^{subformula}]
The !!{subformula}s of $!A$ are $!A$ itself together with all its
proper !!{subformula}s.
\end{defn}

\begin{explain}
Note the subtle difference in how we have defined immediate
!!{subformula}s and proper !!{subformula}s.  In the first case, we
have directly defined the immediate !!{subformula}s of a formula~$!A$
for each possible form of~$!A$.  It is an explicit definition by
cases, and the cases mirror the inductive definition of the set of
!!{formula}s.  In the second case, we have also mirrored the way the
set of all !!{formula}s is defined, but in each case we have also
included the proper !!{subformula}s of the smaller !!{formula}s $!B$,
$!C$ in addition to these !!{formula}s themselves.  This makes the
definition \emph{recursive}.  In general, a definition of a function
on an inductively defined set (in our case, !!{formula}s) is recursive
if the cases in the definition of the function make use of
the function itself. To be well defined, we must make sure, however,
that we only ever use the values of the function for arguments that
come ``before'' the one we are defining---in our case, when defining
``proper !!{subformula}'' for $(!B \ast !C)$ we only use the proper
!!{subformula}s of the ``earlier'' !!{formula}s $!B$ and $!C$.
\end{explain}

\begin{prop}
\ollabel{prop:subfrm-trans}
Suppose $!B$ is a subformula of $!A$ and $!C$ is a subformula of $!B$.
Then $!C$ is a subformula of $!A$. In other words, the subformula
relation is transitive.
\end{prop}

\begin{prob}
Prove \olref[fol][syn][sbf]{prop:subfrm-trans}.
\end{prob}

\begin{prop}
\ollabel{prop:count-subfrms}
Suppose $!A$ is a formula with $n$ connectives and quantifiers.
Then $!A$ has at most $2n+1$ subformulas.
\end{prop}

\begin{prob}
Prove \olref[fol][syn][sbf]{prop:count-subfrms}.
\end{prob}





\olfileid{fol}{syn}{fseq}

\olsection{Formation Sequences}

Defining !!{formula}s via an inductive definition, and the
complementary technique of proving properties of !!{formula}s via
induction, is an elegant and efficient approach. However, it can
also be useful to consider a more bottom-up, step-by-step approach
to the construction of !!{formula}s, which we do here using the
notion of a \emph{formation sequence}.

To show how terms and !!{formula}s can be introduced in this way
without needing to refer to their inductive definitions, we first
introduce the notion of an arbitrary string of symbols drawn from
some language~$\Lang L$.

\begin{defn}[Strings]
\ollabel{defn:string}
Suppose $\Lang L$ is a first-order language. An \emph{$\Lang
L$-string} is a finite sequence of symbols of~$\Lang L$. Where the
language~$\Lang L$ is clearly fixed by the context, we will often
refer to a $\Lang L$-string simply as a \emph{string}.
\end{defn}

\begin{ex}
For any first-order language $\Lang L$, all
$\Lang L$-!!{formula}s are $\Lang L$-strings, but not
conversely. For example, \[)(\Obj v_0\lif\lexists\] is an
$\Lang L$-string but not an $\Lang L$-!!{formula}.
\end{ex}

\begin{defn}[Formation sequences for terms]
\ollabel{defn:fseq-trm}
A finite sequence of $\Lang L$-strings $\tuple{t_0,\dotsc,t_n}$ is a
\emph{formation sequence} for a term $t$ if $t \ident t_n$ and for all
$i \leq n$, either $t_i$ is !!a{variable} or !!a{constant}, or $\Lang
L$ contains a $k$-ary !!{function}~$f$ and there exist
$m_0,\dotsc,m_k < i$ such that $t_i \ident f(t_{m_0},\dotsc,t_{m_k})$.
When it is necessary to distinguish, we will refer to formation
sequences for terms as \emph{term formation sequences}.
\end{defn}

\begin{ex}
The sequence
\[
    \tuple{\Obj c_0, \Obj v_0, \Atom{\Obj f^2_0}{\Obj c_0, \Obj v_0}, \Atom{\Obj f^1_0}{\Atom{\Obj f^2_0}{\Obj c_0, \Obj v_0}}}
\]
is a formation sequence for the term $\Atom{\Obj f^1_0}{\Atom{\Obj
f^2_0}{\Obj c_0, \Obj v_0}}$, as is
\[
    \tuple{\Obj v_0, \Obj c_0, \Atom{\Obj f^2_0}{\Obj c_0, \Obj v_0}, \Atom{\Obj f^1_0}{\Atom{\Obj f^2_0}{\Obj c_0, \Obj v_0}}}.
\]
\end{ex}

\begin{defn}[Formation sequences for formulas]
\ollabel{defn:fseq-frm}
A finite sequence of $\Lang L$-strings $\tuple{!A_0,\dotsc,!A_n}$
is a \emph{formation sequence} for~$!A$ if $!A \ident !A_n$ and
for all $i \leq n$, either $!A_i$ is an atomic !!{formula} or there
exist $j,k < i$ and !!a{variable}~$x$ such that one of the following
holds:
\begin{enumerate}
    $!A_i \ident \lnot !A_j$.
    $!A_i \ident (!A_j \land !A_k)$.
    $!A_i \ident (!A_j \lor !A_k)$.
    $!A_i \ident (!A_j \lif !A_k)$.
    
    $!A_i \ident \lforall[x][!A_j]$.
    $!A_i \ident \lexists[x][!A_j]$.
\end{enumerate}
When it is necessary to distinguish, we will refer to formation
sequences for formulas as \emph{formula formation sequences}.
\end{defn}

\begin{ex}
\[
\tuple{
    \Atom{\Obj A^1_0}{\Obj v_0},
    \Atom{\Obj A^1_1}{\Obj c_1},
    (\Atom{\Obj A^1_1}{\Obj c_1} \land \Atom{\Obj A^1_0}{\Obj v_0}),
    \lexists[\Obj v_0][(\Atom{\Obj A^1_1}{\Obj c_1} \land \Atom{\Obj A^1_0}{\Obj v_0})]
}
\]
is a formation sequence of $\lexists[\Obj v_0][(\Atom{\Obj A^1_1}{\Obj
c_1} \land \Atom{\Obj A^1_0}{\Obj v_0})]$, as is
\begin{multline*}
\tuple{
    \Atom{\Obj A^1_0}{\Obj v_0},
    \Atom{\Obj A^1_1}{\Obj c_1},
    (\Atom{\Obj A^1_1}{\Obj c_1} \land \Atom{\Obj A^1_0}{\Obj v_0}),
    \Atom{\Obj A^1_1}{\Obj c_1},\\
    \lforall[\Obj v_1][\Atom{\Obj A^1_0}{\Obj v_0}],
    \lexists[\Obj v_0][(\Atom{\Obj A^1_1}{\Obj c_1} \land \Atom{\Obj A^1_0}{\Obj v_0})]
}.
\end{multline*}

As can be seen from the second example, formation sequences
may contain ``junk'': !!{formula}s which are redundant or do not
contribute to the construction.
\end{ex}

\begin{prop}\ollabel{prop:formed}
Every !!{formula}~$!A$ in~$\Frm[L]$ has a formation sequence.
\end{prop}

\begin{proof}
Suppose $!A$ is atomic. Then the sequence $\tuple{!A}$ is a
formation sequence for~$!A$.

Now suppose that $!B$ and~$!C$ have formation sequences
$\tuple{!B_0,\dotsc,!B_n}$ and $\tuple{!C_0,\dotsc,!C_m}$
respectively.

\begin{enumerate}
  If $!A \ident \lnot !B$,
      then $\tuple{!B_0,\dotsc,!B_n,\lnot !B_n}$
      is a formation sequence for~$!A$.
  If $!A \ident (!B \land !C)$,
      then $\tuple{!B_0,\dotsc,!B_n,!C_0,\dotsc,!C_m,(!B_n \land !C_m)}$
      is a formation sequence for~$!A$.
  If $!A \ident (!B \lor !C)$,
      then $\tuple{!B_0,\dotsc,!B_n,!C_0,\dotsc,!C_m,(!B_n \lor !C_m)}$
      is a formation sequence for~$!A$.
  If $!A \ident (!B \lif !C)$,
      then $\tuple{!B_0,\dotsc,!B_n,!C_0,\dotsc,!C_m,(!B_n \lif !C_m)}$
      is a formation sequence for~$!A$.
  
  If $!A \ident \lforall[x][!B]$,
      then $\tuple{!B_0,\dotsc,!B_n,\lforall[x][!B_n]}$
      is a formation sequence for~$!A$.
  If $!A \ident \lexists[x][!B]$,
      then $\tuple{!B_0,\dotsc,!B_n,\lexists[x][!B_n]}$
      is a formation sequence for~$!A$.
\end{enumerate}
By the principle of induction on !!{formula}s,
every !!{formula} has a formation sequence.
\end{proof}

We can also prove the converse. This is important because it shows
that our two ways of defining formulas are equivalent: they give
the same results. It also means that we can prove theorems about
formulas by using ordinary induction on the length of formation
sequences.

\begin{lem}
\ollabel{lem:fseq-init}
Suppose that $\tuple{!A_0,\dotsc,!A_n}$ is a formation sequence
for~$!A_n$, and that $k \leq n$. Then $\tuple{!A_0,\dotsc,!A_k}$
is a formation sequence for~$!A_k$.
\end{lem}

\begin{proof}
Exercise.
\end{proof}

\begin{prob}
Prove \olref[fol][syn][fseq]{lem:fseq-init}.
\end{prob}

\begin{thm}
\ollabel{thm:fseq-frm-equiv}
$\Frm[L]$ is the set of all $\Lang L$-strings~$!A$ such that
there exists a formula formation sequence for~$!A$.
\end{thm}

\begin{proof}
Let $F$ be the set of all strings of symbols in the language~$\Lang L$
that have a formation sequence. We have seen in
\olref[fol][syn][fseq]{prop:formed} that $\Frm[L] \subseteq F$, so now
we prove the converse.

Suppose $!A$ has a formation sequence $\tuple{!A_0,\dotsc,!A_n}$.
We prove that $!A \in \Frm[L]$ by strong induction on~$n$.
Our induction hypothesis is that every string of symbols with a
formation sequence of length $m < n$ is in $\Frm[L]$.
By the definition of a formation sequence, either $!A \ident !A_n$ is
atomic or there must exist $j,k < n$ such that one of the
following is the case:
\begin{enumerate}
    $!A \ident \lnot !A_j$.
    $!A \ident (!A_j \land !A_k)$.
    $!A \ident (!A_j \lor !A_k)$.
    $!A \ident (!A_j \lif !A_k)$.
    
    $!A \ident \lforall[x][!A_j]$.
    $!A \ident \lexists[x][!A_j]$.
\end{enumerate}
Now we reason by cases. If $!A$ is atomic then
$!A_n \in \Frm[L_0]$. Suppose instead that
$!A \equiv (!A_j \land !A_k)$. By 
\olref[fol][syn][fseq]{lem:fseq-init},
$\tuple{!A_0,\dotsc,!A_j}$ and $\tuple{!A_0,\dotsc,!A_k}$ are
formation sequences for $!A_j$ and~$!A_k$, respectively. Since
these are proper initial subsequences of the formation sequence
for~$!A$, they both have length less than~$n$. Therefore by
the induction hypothesis, $!A_j$ and~$!A_k$ are in~$\Frm[L_0]$,
and by the definition of !!a{formula}, so is
$(!A_j \land !A_k)$. The other cases follow by parallel
reasoning.
\end{proof}

Formation sequences for terms have similar properties to those
for !!{formula}s.

\begin{prop}
\ollabel{prop:fseq-trm-equiv}
$\Trm[L]$ is the set of all $\Lang L$-strings $t$
such that there exists a term formation sequence for~$t$.
\end{prop}

\begin{proof}
Exercise.
\end{proof}

\begin{prob}
Prove \olref[fol][syn][fseq]{prop:fseq-trm-equiv}.
Hint: use a similar strategy to that used in the proof of
\olref[fol][syn][fseq]{thm:fseq-frm-equiv}.
\end{prob}

There are two types of ``junk'' that can appear in formation
sequences: repeated elements, and elements that are irrelevant
to the construction of the formation or term. We can eliminate
both by looking at minimal formation sequences.

\begin{defn}[Minimal formation sequences]
\ollabel{defn:minimal-fseq}
A formation sequence $\tuple{!A_0, \ldots, !A_n}$ for a
formula~$!A$ is a \emph{minimal formation sequence} for~$!A$
if for every other formation sequence~$s$ for~$!A$,
the length of~$s$ is greater than or equal to~$n+1$.

Similarly, a formation sequence $\tuple{t_0, \ldots, t_n}$
for a term~$t$ is a \emph{minimal formation sequence}
for~$t$ if for every other formation sequence~$s$ for~$t$,
the length of~$s$ is greater than or equal to~$n+1$.
\end{defn}

Note that a formula or term can have more than one minimal
formation sequence, but they will contain exactly the same
strings.

\begin{prop}
\ollabel{prop:subformula-equivs}
The following are equivalent:
\begin{enumerate}
    \item $!B$ is a sub-!!{formula} of~$!A$.
    \item $!B$ occurs in every formation sequence of~$!A$.
    \item $!B$ occurs in a minimal formation sequence of~$!A$.
\end{enumerate}
\end{prop}

\begin{proof}
Exercise.
\end{proof}

\begin{prob}
Prove \olref[fol][syn][fseq]{prop:subformula-equivs}.
\end{prob}

\begin{history}
Formation sequences were introduced by Raymond Smullyan in his
textbook \emph{First-Order Logic} \citep{Smullyan1968}.
Additional properties of formation sequences were established by
\citet{Zuckerman1973}.
\end{history}





\olfileid{fol}{syn}{fvs}

\olsection{Free \printtoken{P}{variable} and \printtoken{P}{sentence}}

\begin{defn}[Free occurrences of !!a{variable}]
\ollabel{defn:free-occ}
The \emph{free} occurrences of !!a{variable} in !!a{formula} are defined
inductively as follows:
\begin{enumerate}
\item \indcase*{!A}{$!A$ is atomic}{all !!{variable} occurrences in
  $\indfrm$ are free.}

\indcase{!A}{\lnot !B}{the free !!{variable}
  occurrences of $\indfrm$ are exactly those of $!B$.}

\item \indcase{!A}{(!B \ast !C)}{the free
  !!{variable} occurrences of $\indfrm$ are those in $!B$
  together with those in~$!C$.}

\indcase{!A}{\lforall[x][!B]}{the free !!{variable}
  occurrences in $\indfrm$ are all of those in~$!B$ except for
  occurrences of~$x$.}

\indcase{!A}{\lexists[x][!B]}{the free !!{variable}
  occurrences in $\indfrm$ are all of those in~$!B$ except for
  occurrences of~$x$.}
\end{enumerate}
\end{defn}

\begin{defn}[Bound Variables]
An occurrence of !!a{variable} in a formula~$!A$ is \emph{bound} if
it is not free.
\end{defn}

\begin{prob}
Give an inductive definition of the bound variable occurrences along
the lines of \olref[fol][syn][fvs]{defn:free-occ}.
\end{prob}

\begin{defn}[Scope]
If $\lforall[x][!B]$ is an occurrence of a subformula
  in a formula~$!A$, then the corresponding occurrence of~$!B$ in~$!A$
  is called the \emph{scope} of the corresponding occurrence
  of~$\lforall[x]$. Similarly for $\lexists[x]$.

If $!B$ is the scope of a quantifier occurrence
$\lforall[x]$ or
    $\lexists[x]$ in~$!A$, then the free occurrences of
$x$ in~$!B$ are bound in $\lforall[x][!B]$ and
$\lexists[x][!B]$. We say that these
occurrences are \emph{bound by} the
mentioned quantifier occurrence.
\end{defn}

\begin{ex}
Consider the following formula:
\[
\lexists[\Obj v_0][\underbrace{\Atom{\Obj A^2_0}{\Obj v_0,\Obj v_1}}_{!B}] 
\]
$!B$ represents the scope of $\lexists[\Obj v_0]$. 
The quantifier binds the occurrence of $\Obj v_0$ in $!B$, but
does not bind the occurrence of $\Obj v_1$. So $\Obj v_1$ is
a free variable in this case.


We can now see how this might work in a more complicated 
!!{formula}~$!A$:
\[
\lforall[\Obj v_0][\underbrace{(\Atom{\Obj A^1_0}{\Obj v_0} \lif
    \Atom{\Obj A^2_0}{\Obj v_0, \Obj v_1})}_{!B}] \lif \lexists[\Obj
  v_1][\underbrace{(\Atom{\Obj A^2_1}{\Obj v_0, \Obj v_1} \lor \lforall[\Obj v_0][\overbrace{\lnot \Atom{\Obj A^1_1}{\Obj v_0}}^{!D}])}_{!C}]
\]
$!B$ is the scope of the first $\lforall[\Obj v_0]$, $!C$ is the scope
of $\lexists[\Obj v_1]$, and $!D$ is the scope of the second
$\lforall[\Obj v_0]$.  The first $\lforall[\Obj v_0]$ binds the
occurrences of $\Obj v_0$ in~$!B$, $\lexists[\Obj v_1]$ binds the occurrence
of $\Obj v_1$ in $!C$, and the second $\lforall[\Obj v_0]$ binds the
occurrence of $\Obj v_0$ in~$!D$.  The first occurrence of $\Obj v_1$
and the fourth occurrence of $\Obj v_0$ are free in~$!A$. The last
occurrence of $\Obj v_0$ is free in $!D$, but bound in $!C$ and~$!A$.
\end{ex}

\begin{defn}[Sentence]
!!^a{formula}~$!A$ is \article{sentence} \emph{!!{sentence}} iff it
contains no free occurrences of !!{variable}s.
\end{defn}







\olfileid{fol}{syn}{sub}

\olsection{Substitution}

\begin{defn}[Substitution in a term]
We define $\Subst{s}{t}{x}$, the result of \emph{substituting} $t$
for every occurrence of~$x$ in $s$, recursively:
\begin{enumerate}
\item \indcase{s}{c}{$\Subst{\indfrm}{t}{x}$ is just $s$.}

\item \indcase{s}{y}{$\Subst{\indfrm}{t}{x}$ is also just~$s$,
  provided $y$ is a variable and $y \not\ident x$.}

\item \indcase{s}{x}{$\Subst{\indfrm}{t}{x}$ is~$t$.}

\item\indcase{s}{\Atom{f}{t_1, \dots, t_n}}{$\Subst{\indfrmp}{t}{x}$ is
  $\Atom{f}{\Subst{t_1}{t}{x}, \dots, \Subst{t_n}{t}{x}}$.}
\end{enumerate}
\end{defn}

\begin{defn}
A term~$t$ is \emph{!!{free for}} $x$ in $!A$ if none of the free
occurrences of~$x$ in $!A$ occur in the scope of a quantifier that
binds a variable in~$t$.
\end{defn}

\begin{ex} ~
\begin{enumerate}
\item $\Obj v_8$ is free for $\Obj v_1$ in $\lexists[\Obj
  v_3]\Atom{\Obj A^2_4}{\Obj v_3,\Obj v_1}$

\item $\Obj f^2_1(\Obj v_1, \Obj v_2)$ is \emph{not} free for $\Obj
  v_0$ in $\lforall[\Obj v_2]\Atom{\Obj A^2_4}{\Obj v_0,\Obj v_2}$
\end{enumerate}
\end{ex}

\begin{defn}[Substitution in !!a{formula}]
If $!A$ is !!a{formula}, $x$~is !!a{variable}, and $t$~is a term
!!{free for}~$x$ in~$!A$, then $\Subst{!A}{t}{x}$ is the result of
substituting $t$ for all free occurrences of~$x$ in~$!A$.
\begin{enumerate}
\indcase{!A}{\lfalse}{$\Subst{\indfrm}{t}{x}$ is
    $\lfalse$.}



\item \indcase{!A}{\Atom{P}{t_1,\dots,
    t_n}}{$\Subst{\indfrm}{t}{x}$ is $\Atom{P}{\Subst{t_1}{t}{x},
    \dots, \Subst{t_n}{t}{x}}$.}

\item \indcase{!A}{\eq[t_1][t_2]}{$\Subst{\indfrmp}{t}{x}$ is
  $\Subst{t_1}{t}{x} = \Subst{t_2}{t}{x}$.}

\indcase{!A}{\lnot !B}{$\Subst{\indfrmp}{t}{x}$ is
    $\lnot \Subst{!B}{t}{x}$.}

\indcase{!A}{(!B \land
    !C)}{$\Subst{\indfrmp}{t}{x}$ is $(\Subst{!B}{t}{x} \land
    \Subst{!C}{t}{x})$.}

\indcase{!A}{(!B \lor
    !C)}{$\Subst{\indfrmp}{t}{x}$ is $(\Subst{!B}{t}{x} \lor
    \Subst{!C}{t}{x})$.}

\indcase{!A}{(!B \lif
    !C)}{$\Subst{\indfrmp}{t}{x}$ is $(\Subst{!B}{t}{x} \lif
    \Subst{!C}{t}{x})$.}




\indcase{!A}{\lforall[y][!B]}{$\Subst{\indfrmp}{t}{x}$
    is $\lforall[y][\Subst{!B}{t}{x}]$, provided $y$ is a variable
    other than $x$; otherwise $\Subst{\indfrmp}{t}{x}$
    is just $\indfrm$.}


\indcase{!A}{\lexists[y][!B]}{$\Subst{\indfrmp}{t}{x}$
    is $\lexists[y][\Subst{!B}{t}{x}]$, provided $y$ is a variable
    other than $x$; otherwise $\Subst{\indfrmp}{t}{x}$
    is just $\indfrm$.}
\end{enumerate}
\end{defn}

\begin{explain}
Note that substitution may be vacuous: If $x$ does not occur in $!A$
at all, then $\Subst{!A}{t}{x}$ is just~$!A$.

The restriction that $t$ must be !!{free for}~$x$ in~$!A$ is necessary to
exclude cases like the following. If $!A \ident \lexists[y][x < y]$
and $t \ident y$, then $\Subst{!A}{t}{x}$ would be $\lexists[y][y <
  y]$. In this case the free variable $y$ is ``captured'' by the
quantifier $\lexists[y]$ upon substitution, and that is undesirable.
For instance, we would like it to be the case that whenever
$\lforall[x][!B]$ holds, so does $\Subst{!B}{t}{x}$. But consider
$\lforall[x][\lexists[y][x < y]]$ (here $!B$ is $\lexists[y][x <
  y]$). It is a sentence that is true about, e.g., the natural numbers:
for every number~$x$ there is a number~$y$ greater than it. If we
allowed $y$ as a possible substitution for~$x$, we would end up with
$\Subst{!B}{y}{x} \ident \lexists[y][y < y]$, which is false. We
prevent this by requiring that none of the free variables in~$t$ would
end up being bound by a quantifier in~$!A$.
\end{explain}

We often use the following convention to avoid cumbersome notation: If
$!A$ is !!a{formula} which may contain the !!{variable}~$x$ free, we
also write~$!A(x)$ to indicate this. When it is clear which $!A$
and~$x$ we have in mind, and $t$ is a term (assumed to be free for $x$
in $!A(x)$), then we write $!A(t)$ as short for $\Subst{!A}{t}{x}$. So
for instance, we might say, ``we call $!A(t)$ an instance
of~$\lforall[x][!A(x)]$.'' By this we mean that if $!A$~is any
!!{formula}, $x$~!!a{variable}, and $t$~a term that's free for~$x$
in~$!A$, then $\Subst{!A}{t}{x}$ is an instance of~$\lforall[x][!A]$.



\OLEndChapterHook





